\section{Introduction}

The overarching goal of this project is to develop a domain-specific, pan-cancer single-cell foundation
model capable of predicting pharmacological response, both efficacy and toxicity, that can subsequently be
fine-tuned for specific cancer types or for clinical datasets with binary outcomes. Tumor heterogeneity is
a major driver of drug resistance and treatment failure, yet this resolution is obscured whenever drug
response is modeled from bulk RNA-seq. Predicting response at the single-cell level would allow a model to
report a distribution of sensitivities within a sample and, in principle, to identify rare, naturally
resistant subclones before selection pressure acts on them. The immediate obstacle is that single-cell
drug-screening data does not exist at scale.

To make progress despite this, the project adopts a weakly supervised strategy in which bulk
pharmacological labels are mapped onto the individual cells of the corresponding \emph{in vitro} cell
lines, and the model is trained to map heterogeneous single-cell inputs to the average bulk response.
Laboratory cell lines do not fully recapitulate clinical tumors, but they provide the high-volume,
continuous labels required to train a robust representation that can later be fine-tuned on scarcer
clinical data. The central methodological difficulty in this setup is a resolution mismatch: standard
pharmacological labels, including those of CTRPv2 \cite{ctrp}, PRISM \cite{prism} and GDSC \cite{gdsc}, are
derived from bulk cell populations, whereas the transcriptomic data from SCP542 \cite{scp542} captures
individual cellular heterogeneity. Mapping one average bulk score onto hundreds of heterogeneous cells
introduces substantial label noise, and regressing raw, sparse single-cell counts directly onto these
labels is unstable.

This motivates the use of a foundation-model representation as an intermediate prior. An exploratory
latent-space analysis illustrates the effect (Figure~\ref{fig:umap}). When cells are colored by cancer
type, full-transcriptome PCA is dominated by tissue-specific markers and separates the data into distinct
lineage clusters, so that a downstream model could achieve low error simply by classifying the cell line.
The scGPT embeddings \cite{scgpt} instead project the same cells onto a continuous, shared manifold in
which tissues are intermixed. The project therefore diverges from prior work such as PERCEPTION
\cite{perception}, which pre-trained on bulk data to predict PRISM response, by focusing on reusable
representation learning: it evaluates whether transformer-based single-cell embeddings improve
drug-response prediction over a standard PCA baseline, and it is designed to extend to multi-task learning
across several databases. The central hypothesis is that scGPT embeddings act as a denoised biological
prior that aligns functional cellular states across tissues, so that a regressor trained on this
representation is encouraged to learn transcriptomic drivers of sensitivity rather than to memorize the
tissue of origin.

This report covers only the initial phase of the work: a single database (CTRPv2), a single continuous
response score, single-cell inputs and a supervised regressor. The multi-task run reported here spans
\NDrugs{} drugs but a single database and metric; cross-database integration, toxicity and clinical
fine-tuning are left to later work.

\begin{figure}[t]
\centering
\includegraphics[width=0.9\textwidth]{fig_umap.png}
\caption{\textbf{scGPT removes the tissue-of-origin bias that dominates a PCA baseline.} Cells are embedded
by full-transcriptome PCA (left) and by scGPT (right), reduced to two dimensions by UMAP and colored by
cancer type. PCA separates the data into tissue-specific clusters, whereas scGPT projects the cells onto a
continuous, shared manifold. Produced by \code{notebooks/06\_verify\_variants.ipynb}.}
\label{fig:umap}
\end{figure}
